% a06_BALtorus_A3: continuation of a06_BALtorus_A2 (input from it). Group (A), part 3 of 4:
% the two-sided depth factor of Ba4's Theorem, the domain differences (5.5) and (1.20), Ba1 Prop. 2.2, the summary theorem, the polylog bound.

The remaining estimates use the printed theorem of \cite{Ba4}, which is \cite[Prop.~2.1]{Ba1}: for $\alpha<1$ there are $\delta_0,c_0,R_0>0$, independent of $A$, $k$, $\Omega$, depending on $d$, $M$ (and $c_0$ on $\alpha$), such that for $e_k$ small and $f\colon\Omega\to\R^N$: the $C^{1,\alpha}$ bound \cite[(1.9)]{Ba4} for $x,x'\in\Omega$ with $\dist(\{x,x'\},\Omega^c)\ge R_0$; $|(D^\eta_{A,\mu}G_k(\Omega,A)f)(x)|,\,|(G_k(\Omega,A)f)(x)|\le c_0e^{-\delta_0\dist(x,\supp f)}\norm f_\infty$ for $x\in\Omega$ with $\dist(x,\Omega^c)\ge R_0$ \cite[(1.10)]{Ba4}; and, for $\Omega\subset\Omega_0$ and $\delta G_k(\Omega,\Omega_0,A)\defeq G_k(\Omega,A)-G_k(\Omega_0,A)$ \cite[(1.11)]{Ba4}, the same bounds with the same restrictions and an additional factor \cite[(1.12)]{Ba4}, printed in \cite[(2.26)]{Ba1} as $\exp(-\delta_0\dist(\supp f,\Omega^c)-\delta_0\dist(\{x,x'\},\Omega^c))$. When $\Omega$ is the whole torus the depth condition is void \cite[p.~611]{Ba1}. Only this restricted form (points at depth at least $R_0$, or $\Omega$ the whole torus) is used anywhere in this appendix; the clause of \cite[p.~611]{Ba1} and \cite[p.~573]{Ba4} that drops the depth condition on rectangular parallelepipeds is not used.

\begin{lemma}[the depth factor, two-sided]\label{lem:a06A:twosided}
Let $\Omega\subset\Omega_0$ be unions of large blocks, $A$ regular on $\Omega_0$, $e_k$ small, $\supp f$ of diameter at most $s$, and $x\in\Omega$ with $\dist(x,\Omega^c)\ge R_0$. Then
\[
|(\delta G_k(\Omega,\Omega_0,A)f)(x)|\le c_0e^{\delta_0s/2}\exp\big[-\tfrac{\delta_0}2\big(\dist(x,\supp f)+\dist(x,\Omega^c)+\dist(\supp f,\Omega^c)\big)\big]\norm f_\infty,
\]
and the same for $D^\eta_{A,\mu}\delta G_k(\Omega,\Omega_0,A)f$, and for the H\"older quotient of \cite[(1.9)]{Ba4} with $|x-x'|\le1$ and $\dist(\{x,x'\},\Omega^c)\ge R_0$, with $\dist(\{x,x'\},\cdot)$ in place of $\dist(x,\cdot)$. For general $f$ (no bound on the diameter) the bounds hold at the rate $\delta_0/4$ with the constant $c_0e^{3\delta_0/4}S(\delta_0/4)$ ($c_0e^{\delta_0}S(\delta_0/4)$ for the H\"older quotient).
\end{lemma}

\begin{proof}
$\dist(x,\Omega^c)\le\dist(x,\supp f)+s+\dist(\supp f,\Omega^c)$, so $\dist(x,\supp f)+\dist(x,\Omega^c)+\dist(\supp f,\Omega^c)\le2\dist(x,\supp f)+2\dist(\supp f,\Omega^c)+s$. The printed factor, read with both of its terms as $\dist(\supp f,\Omega^c)$, gives $c_0e^{-\delta_0\dist(x,\supp f)}e^{-2\delta_0\dist(\supp f,\Omega^c)}\norm f_\infty$, which is at most the right side; read as printed in \cite[(2.26)]{Ba1}, it gives the right side without the factor $e^{\delta_0s/2}$ and at the rate $\delta_0$. So the bound holds under either reading. For the H\"older quotient the printed factor carries $\dist(\{x,x'\},\cdot)$ and the same arithmetic applies, because $\dist(\{x,x'\},\Omega^c)\le\dist(\{x,x'\},\supp f)+s+\dist(\supp f,\Omega^c)$ (the inequality for $x$ and for $x'$ separately, then the minimum over the two points). For general $f$ decompose $\supp f$ into the unit blocks $\square$ it meets, apply the bound to $f\mathbf 1_\square$ ($s\le1$; $\dist(x,\square\cap\supp f)\ge\dist(x,\supp f)$ and $\dist(\square\cap\supp f,\Omega^c)\ge\dist(\supp f,\Omega^c)$), and sum: with $\square_n$ the unit block of $n\in\Z^d$ and $x\in\square_{n_x}$, $\dist(x,\square_n)\ge(|n-n_x|-1)_+$, so $\sum_\square e^{-(\delta_0/2)\dist(x,\square\cap\supp f)}\le e^{-(\delta_0/4)\dist(x,\supp f)}\sum_{n\in\Z^d}e^{-(\delta_0/4)(|n-n_x|-1)_+}\le e^{-(\delta_0/4)\dist(x,\supp f)}e^{\delta_0/4}S(\delta_0/4)$; for the H\"older quotient $\dist(\{x,x'\},\square_n)\ge\dist(x,\square_n)-1$ costs one more factor $e^{\delta_0/4}$.
\end{proof}

\begin{lemma}[{domain differences; \cite[(5.5), (1.20)]{Ba4}}]\label{lem:a06A:domdiff}
Let $\Omega\subset\Omega_0$ be unions of large blocks, (H1)--(H3), $A$ regular on $\Omega_0$, $e_k\le e_1$ and small enough for the printed theorem. Put
\[
\delta_7\defeq\min\{\delta_0/2,\,\theta_0(\gamma_\emptyset,1)\},\qquad\delta^S_7\defeq\min\{\theta_0,\,\delta_7\}=\min\{\theta_0(\gamma_L,L),\,\delta_0/2\}.
\]
\begin{enumerate}[label=\textup{(\alph*)},nosep,leftmargin=*]
\item For $y,y'\in\Omega^{(k)}$, with $d(y)\defeq\dist(B^k(y),\Omega^c)$,
\[
|(\Delta^{(k)}(\Omega,A)-\Delta^{(k)}(\Omega_0,A))(y,y')|\le c_6\exp\big[-\delta_7\big(|y-y'|+d(y)+d(y')\big)\big],
\]
$c_6\defeq\max\{a^2c_0e^{\delta_0/2},\,4a^2\gamma_\emptyset^{-1}e^{2\theta_0(\gamma_\emptyset,1)R_0}\}e^{\delta_7}$: this is \cite[(5.5)]{Ba4}.
\item For every $\Lambda\subset\Omega^{(k)}$ that is a union of $L$-blocks and $x,x'\in\Lambda$, with $\tilde d(x)\defeq\dist(x,\Omega^{(k)c})$,
\[
|C_\Lambda(\Omega,A;x,x')-C_\Lambda(\Omega_0,A;x,x')|\le c_{20}\exp\big[-\tfrac{\delta^S_7}4\big(|x-x'|+\tilde d(x)+\tilde d(x')\big)\big],
\]
$c_{20}\defeq c_2^2c_6e^{\delta^S_7/2}S(\delta^S_7/2)^2$: this is \cite[(1.20)]{Ba4}.
\end{enumerate}
\end{lemma}

\begin{proof}
(a) $Q_k(A)$ on $\Omega$ and on $\Omega_0$ agree on the blocks $B^k(y)$, $y\in\Omega^{(k)}$, so by \cite[(1.14)]{Ba4} the kernel is $\delta\Delta(y,y')=-a_k^2\langle q_k(y),\delta G_k(\Omega,\Omega_0,A)q_k(y')\rangle$, with $q_k(y),q_k(y')$ supported in $\Omega$. Three cases. If $d(y)\ge R_0$, every $x\in B^k(y)$ has $\dist(x,\Omega^c)\ge R_0$, and \cref{lem:a06A:twosided} with $f=q_k(y')$ ($\norm f_\infty=1$, $s\le1$) gives $|(\delta G_kq_k(y'))(x)|\le c_0e^{\delta_0/2}e^{-(\delta_0/2)(\dist(B^k(y),B^k(y'))+d(y)+d(y'))}$ for $x\in B^k(y)$; pairing with $q_k(y)$, whose entries have modulus $1$ on a block of volume $1$, gives $|\delta\Delta(y,y')|\le a_k^2c_0e^{\delta_0/2}e^{-(\delta_0/2)(\dist(B^k(y),B^k(y'))+d(y)+d(y'))}$. If $d(y')\ge R_0$, use $\delta\Delta(y,y')=\delta\Delta(y',y)^{\mathsf T}$. If $d(y),d(y')<R_0$, apply \cref{lem:a06A:CT}(i) to $G_k(\Omega,A)$ and to $G_k(\Omega_0,A)$ separately (\cref{lem:a06A:coer} with $\Lambda=\emptyset$ on $\Omega$ and on $\Omega_0$; $\ell_*=1$, $\gamma=\gamma_\emptyset$): $|\delta\Delta(y,y')|\le4a_k^2\gamma_\emptyset^{-1}e^{-\theta_0(\gamma_\emptyset,1)\dist(B^k(y),B^k(y'))}\le4a_k^2\gamma_\emptyset^{-1}e^{2\theta_0(\gamma_\emptyset,1)R_0}e^{-\theta_0(\gamma_\emptyset,1)(\dist(B^k(y),B^k(y'))+d(y)+d(y'))}$. Finally $\dist(B^k(y),B^k(y'))\ge|y-y'|-1$.

(b) The inverses are bounded by \cref{lem:a06A:blockinv}(b) on both domains, so $C_\Lambda(\Omega,A)-C_\Lambda(\Omega_0,A)=X_\Lambda(\Omega,A)^{-1}[X_\Lambda(\Omega_0,A)-X_\Lambda(\Omega,A)]X_\Lambda(\Omega_0,A)^{-1}$, and $X_\Lambda(\Omega_0,A)-X_\Lambda(\Omega,A)=-\delta\Delta|_\Lambda$, because $\Lambda\subset\Omega^{(k)}$ and $P(A)$ is the same on $\Lambda$ for both domains. With \cref{lem:a06A:blockinv}(c) for both domains and (a),
\[
|\delta C(x,x')|\le c_2^2c_6\sum_{y,y'\in\Lambda}e^{-\delta^S_7\Sigma},\qquad\Sigma\defeq|x-y|+|y-y'|+d(y)+d(y')+|y'-x'| .
\]
Since $\Omega^c=B^k(\Omega^{(k)c})$ and $\dist(B^k(y),B^k(z))\ge|y-z|-1$, $\dist(y,\Omega^{(k)c})\le d(y)+1$; so $\tilde d(x)\le|x-y|+d(y)+1$, $\tilde d(x')\le|x'-y'|+d(y')+1$, and $|x-x'|+\tilde d(x)+\tilde d(x')\le2\Sigma+2$. Hence $e^{-\delta^S_7\Sigma}\le e^{\delta^S_7/2}e^{-(\delta^S_7/4)(|x-x'|+\tilde d(x)+\tilde d(x'))}e^{-(\delta^S_7/2)(|x-y|+|y'-x'|)}$, and the sums over $y$ and $y'$ give $S(\delta^S_7/2)^2$.
\end{proof}

\begin{proposition}[{\cite[Prop.~2.2, (2.27)--(2.29)]{Ba1}}]\label{prop:a06A:Ba1P22}
Assume (H1)--(H3), $A$ regular on $\Omega_0\supset\Omega$, $e_k$ as in \cref{lem:a06A:domdiff}. Then for $y,y'\in\Omega^{(k)}$
\[
|\Delta^{(k)}(\Omega,A;y,y')|\le c_7e^{-\delta_7|y-y'|},\qquad|\Delta^{(k)}(\Omega,A;y,y')-\Delta^{(k)}(\Omega_0,A;y,y')|\le c_7\exp\big[-\delta_7\big(|y-y'|+\dist(y,\Omega^{(k)c})+\dist(y',\Omega^{(k)c})\big)\big],
\]
with $c_7\defeq\max\{a+2a^2\gamma_\emptyset^{-1}e^{\theta_0(\gamma_\emptyset,1)},\,c_6e^{2\delta_7}\}$. The first bound holds on every union $\Omega$ of unit blocks and uses only \cref{lem:a06A:coer,lem:a06A:CT}; its rate is in fact $\theta_0(\gamma_\emptyset,1)\ge\delta_7$.
\end{proposition}

\begin{proof}
$\Delta^{(k)}(y,y')=a_k\delta_{yy'}-a_k^2\langle q_k(y),G_k(\Omega,A)q_k(y')\rangle$ \cite[(1.14)]{Ba4}. \Cref{lem:a06A:coer} with $\Lambda=\emptyset$ and \cref{lem:a06A:CT}(i) with $\ell_*=1$ give $|\langle q_k(y),G_kq_k(y')\rangle|\le2\gamma_\emptyset^{-1}e^{-\theta_0(\gamma_\emptyset,1)\dist(B^k(y),B^k(y'))}$, and $\dist(B^k(y),B^k(y'))\ge|y-y'|-1$; this is the first bound. The second is \cref{lem:a06A:domdiff}(a) with $d(y)\ge\dist(y,\Omega^{(k)c})-1$.
\end{proof}

\begin{theorem}[the propagator estimates of group (A)]\label{thm:a06A:groupA}
Let $d=3$, $L=5$, (H1)--(H3), $\Omega\subset\Omega_0$ unions of large blocks of $T^{(k)}_1$ (the whole torus included), $\Lambda\subset\Omega^{(k)}$ a union of $L$-blocks, $A$ regular on $\Omega_0$, and $\eps_0$ below the caps of \cref{lem:a06A:cube,lem:a06A:domdiff}. Then \cite[(1.8)]{Ba4} on $\Omega$ (\cref{lem:a06A:coer}); the $G_k$ clause of \cite[Cor.~2.3]{Ba4} for $G_k(\Omega,A)$, $G_k(\Omega,\Lambda,A)$ and $G^\eta_{k+1}(\Omega,A)$ (\cref{lem:a06A:CT}); \cite[Prop.~2.3]{Ba1}, i.e.\ \cite[(1.15), (1.16), (1.18), (1.20)]{Ba4} (\cref{lem:a06A:blockinv}(b), (c), \cref{lem:a06A:CLdiff}, \cref{lem:a06A:domdiff}(b)); \cite[(5.5)]{Ba4} (\cref{lem:a06A:domdiff}(a)); and \cite[Prop.~2.2]{Ba1} (\cref{prop:a06A:Ba1P22}) hold with the constants
\[
\gamma_\emptyset,\ \gamma_L,\ \gamma_X,\ \gamma_1,\ \theta_0=\theta_0(\gamma_L,L),\ \theta_0(\gamma_\emptyset,1),\ c_2,\ c_3,\ c_4,\ c_6,\ c_7,\ c_{20},\ \delta_7,\ \delta^S_7
\]
defined above. They depend on $(d,N,\norm q,c,\beta,a,L)$, on $M$ through the constants $\delta_0,c_0,R_0$ of \cite[Thm.]{Ba4} (which enter \cref{lem:a06A:twosided,lem:a06A:domdiff}, \cref{prop:a06A:Ba1P22} and nothing else), and on the cap on $\eps_0$; they do not depend on $\eta$, $k$, $\Omega$, $\Omega_0$, $\Lambda$, $A$, $\eps$ or the torus.
\end{theorem}

\begin{proof}
Each estimate is the lemma named after it; the dependences are read off the definitions of the constants. The order of the arguments is not circular: \cref{lem:a06A:cube,lem:a06A:coer,lem:a06A:CT,lem:a06A:blockinv,lem:a06A:resdiff,lem:a06A:CLdiff} and the first bound of \cref{prop:a06A:Ba1P22} use no estimate of \cite{Ba4} beyond its definitions, while \cref{lem:a06A:twosided,lem:a06A:domdiff} and the second bound of \cref{prop:a06A:Ba1P22} use the printed theorem of \cite{Ba4} in its restricted form (supplied on a torus, with the constants of \cite{Ba4}, by \cref{lem:a06A:torus} below), whose proof \cite[Sect.~2]{Ba4} rests on \cite[Lemmas~2.1, 2.2, 2.4]{Ba4}; the $j$-uniform form of \cite[Lemma~2.4]{Ba4} is \cref{thm:a06:A5}, whose clause \cite[(2.37)]{Ba4} is \cref{lem:a06A:blockinv}(c) at $A=0$, $\Lambda=\Omega^{(k)}$.
\end{proof}

\begin{remark}[dependence on $L$; the depth margin]\label{rem:a06A:L}
(i) $\gamma_\emptyset=\frac14\min\{2,a_*\}$ and the constants of \cref{thm:a06:A5} do not depend on $L$, with $a_*\ge3a/4$ for $L\ge2$; $\gamma_L$, $\gamma_X$, $\theta_0$ and $c_2$ do, through $\gamma_L\propto L^{-2}$, as \cite[p.~588]{Ba4} says of the constants of \cite[(2.37)]{Ba4} and \cite[Prop.~2.3]{Ba1}. At the fixed value $L=5$ all are numbers. (ii) The depth condition of the printed theorem is met wherever it is invoked in Steps~3(B)--3(D): with $r(t)=R(1+\ln t^{-1})^\varrho$ and $R>R_0$ \cite[(2.7)]{Ba2}, every point within distance $1$ of a set at distance more than $2r(t)$ from a boundary has depth more than $2R-1\ge R_0$ (as $R_0\ge1$ may be assumed, enlarging $R_0$ weakening nothing in the theorem); pairs $x,x'$ with $|x-x'|>1$ in \cite[(1.9)]{Ba4} reduce to the sup bounds \cite[p.~578]{Ba4}.
\end{remark}

\begin{lemma}[the polylogarithmic layer]\label{lem:a06A:polylog}
Let $r(t)=R(1+\ln t^{-1})^\varrho$ with $\varrho>1$, and $\delta,R,\kappa>0$. Then $e^{-\delta r(t)}\le t^\kappa$ for all $t\in(0,1]$ if $\delta R\ge\kappa$, and in general for $t\le t_0\defeq\exp(-(\kappa/(\delta R))^{1/(\varrho-1)})$. Consequently $e^{-\delta r(t)}\le C_\kappa t^\kappa$ on $(0,1]$ with $C_\kappa\defeq t_0^{-\kappa}=\exp(\kappa(\kappa/(\delta R))^{1/(\varrho-1)})$.
\end{lemma}

\begin{proof}
With $u\defeq\ln t^{-1}\ge0$, $\delta R(1+u)^\varrho\ge\delta R\max\{1+u,u^\varrho\}$. This is at least $\kappa u$ if $\delta R\ge\kappa$, or if $u^{\varrho-1}\ge\kappa/(\delta R)$, i.e.\ $t\le t_0$. For $t\in[t_0,1]$, $e^{-\delta r(t)}\le1\le t_0^{-\kappa}t^\kappa$.
\end{proof}

% a06_BALtorus_A4: continuation of a06_BALtorus_A3 (input from it). Group (A), part 4 of 4: the torus.

\begin{lemma}[{\cite[Prop.~2.1]{Ba1} on a torus, with torus-free constants}]\label{lem:a06A:torus}
Let $T$ be a torus $\eta\Z^d/\Gamma$ whose period lattice $\Gamma$ is contained in $M\Z^d$ (so the large blocks of the unit lattice tile $T$) and whose sup-systole is at least $3M$; let $\Omega\subset T$ be a union of large blocks ($\Omega=T$ included), $A$ regular on $\Omega$, and $e_k$ small. Then the theorem of \cite{Ba4} (\cite[Prop.~2.1, (2.24)--(2.26)]{Ba1}) holds on $T$ in its restricted form, with distances the quotient sup-distances and with the constants $\delta_0$, $c_0$, $R_0$ of \cite{Ba4}, which do not depend on $\Gamma$; for $\Omega=T$ it holds for all $x,x'\in T$. The same is true for tori of sup-systole less than $3M$ (but at least $M$) whose period lattice lies in $M\Z^d$. Consequently every estimate of \cref{thm:a06A:groupA} holds on $T$ with the constants of \cref{thm:a06A:groupA}.
\end{lemma}

\begin{proof}
The proof of \cite[Sect.~2]{Ba4} (pp.~574--580) is carried onto $T$ step by step; it uses the geometry of the lattice in the following places, and nowhere else.

(1) \emph{The local cubes.} Let $J$ index the vertices $Mj$ of the large-block partition of $T$ and $C_j\defeq Mj+[-M,M]^d$ (in unit-lattice units). Since $2M$ is less than the sup-systole, the projection $\eta\Z^d\to T$ is injective on each $C_j$; it preserves bonds, unit blocks, $L$-blocks and the contours of $Q_k$ inside blocks. So $\Pi_j\defeq\Omega\cap C_j$ \cite[p.~575]{Ba4} is a union of at most $2^d$ large blocks inside an embedded cube of side $2M$, a cube when $Mj$ is not a boundary point of $\Omega$, and the decomposition $A=A_0+A'$ on $C_j$ with the constant $A_0\defeq A(Mj)$ and the local field $A_j\defeq A_0+\theta_jA'$ of \cite[p.~575]{Ba4} is built from $A|_{C_j}$ alone, with $|A'|,|\partial^\eta A'|\le c'e_k^{\beta-1}$, $c'=dMc$. Hence \cite[Lemma~2.1]{Ba4} applies to every $G_k(\Pi_j,A_j)$ and \cite[Lemma~2.2]{Ba4} to the interior cubes (for $\Pi_j$ meeting $\partial\Omega$, $A_j\defeq A$ \cite[p.~575]{Ba4}), with their constants, which depend on $d$, $M$, $\alpha$, $c$ only; \cite[Lemma~2.4]{Ba4}, used in the proof of \cite[Lemma~2.2]{Ba4}, is \cref{thm:a06:A5} on parallelepipeds of the fine lattice.

(2) \emph{The partition of unity.} With $h\in C_0^\infty((-\frac23,\frac23))$, $h=1$ on $[-\frac13,\frac13]$, $\sum_{n\in\Z}h(\cdot-n)^2=1$, and $h_j(x)\defeq\prod_\mu h(x_\mu/M-j_\mu)$ \cite[p.~575]{Ba4}: because $\Gamma\subset M\Z^d$ the family $\{h_j\}_{j\in J}$ is well defined on $T$, $\sum_{j\in J}h_j^2=1$ on $T$ and $\supp h_j\subset C_j$. So the approximate inverse $G_0=\sum_jh_jG_k(\Pi_j,A_j)h_j$ \cite[(2.2)]{Ba4} and the identities \cite[(2.3)--(2.9)]{Ba4} hold on $T$ as on $\eta\Z^d$; for $\Omega=T$ the boundary of $\Omega$ is empty and the boundary term of \cite[(2.6)]{Ba4} vanishes.

(3) \emph{The remainder.} Each $j$ has at most $3^d-1$ neighbours and each point lies in at most $2^d$ supports, so the operator $R$ of \cite[(2.11)]{Ba4} has $\norm R_{2,2}\le c_dO(M^{-1})<1$ by \cite[Lemma~2.1]{Ba4}, as on \cite[p.~577]{Ba4}; the Neumann series \cite[(2.12)]{Ba4} converges and gives the walk representation \cite[(2.13)]{Ba4}, the walks being sequences of indices $j\in J$ with consecutive terms adjacent, and the count $2^d(3^d)^{n-1}2^d$ of walks of length $n$ \cite[p.~579]{Ba4} is unchanged.

(4) \emph{The metric.} Consecutive cubes of a walk have centres at quotient distance at most $M$, so a walk of length $n$ from a cube containing $x,x'$ to a cube containing points of $\supp f$ has $\dist(\{x,x'\},\supp f)\le(n+2)M$, i.e.\ $n\ge M^{-1}\dist(\{x,x'\},\supp f)-2$, which is the restriction used in \cite[(2.22)]{Ba4}; with (3), the geometric sum \cite[(2.22)]{Ba4} gives $\delta_0$ and $c_0$ as on $\eta\Z^d$, once $M$ is fixed with $3^dc_{2.1}O(1)M^{-1}\le\mathrm e^{-1}$, $c_{2.1}$ the constant of \cite[Lemma~2.1, (2.15)]{Ba4} (not the $c_2$ of \cref{lem:a06A:blockinv}). The definition \cite[(2.19)]{Ba4} of $R_0$ (the diameter of the union of the first $n_0$ cubes of a walk, plus $2M$) and the implication ``$\dist(\{x,x'\},\Omega^c)\ge R_0$, hence the first $n_0$ cubes $\Pi_{\omega_i}$ of every walk from $x,x'$ are cubes contained in $\Omega$'' use only the quotient metric and (1); on these cubes $A_j$ is constant near $\partial\Pi_j$ and \cite[Lemma~2.2]{Ba4} applies, on the others only \cite[Lemma~2.1]{Ba4} is used, as in \cite[(2.20)--(2.21)]{Ba4}. The cancellation of the terms with interior cubes in the two representations of $\delta G_k(\Omega,\Omega_0,A)$ \cite[p.~579]{Ba4} uses (1)--(2) only. If $\Omega=T$, every $\Pi_j=C_j$ is a cube with $A_j$ constant near its boundary, so no restriction on $x,x'$ is needed \cite[p.~611]{Ba1}.

(5) \emph{Small tori.} If the sup-systole of $T=\eta\Z^d/\Gamma$ is less than $3M$, let $\tilde T\defeq\eta\Z^d/3\Gamma$, a covering torus of systole at least $3M$ with deck group $(\Z/3)^d$. The operators $-\Delta^{\eta,N}_{A,\Omega}+m_k^2+a_kP_k(A)$ lifted to $\tilde T$ commute with the deck translations and map deck-invariant functions to deck-invariant functions, so for $f$ on $\Omega$ the lift of $G_k(\Omega,A)f$ is $G_k(\tilde\Omega,\tilde A)\tilde f$, with $\tilde\Omega,\tilde A,\tilde f$ the lifts; $\norm{\tilde f}_\infty=\norm f_\infty$, and the quotient distances in $\tilde T$ from a lift of $x$ to $\supp\tilde f$ and to $\tilde\Omega^c$ equal those in $T$. So the bounds on $\tilde T$, proved in (1)--(4), descend to $T$ with the same constants. The same applies to $\delta G_k$.

No other constant of \cite[Sect.~2]{Ba4} involves the global geometry. Finally, \cref{lem:a06A:cube,lem:a06A:coer,lem:a06A:CT,lem:a06A:blockinv,lem:a06A:resdiff,lem:a06A:CLdiff}, the first bound of \cref{prop:a06A:Ba1P22} and \cref{lem:a06A:polylog} are proved above on an arbitrary union of blocks of an arbitrary torus (their only global ingredients are the sums $S(\theta)$ and $S_L(\theta)$ over the block lattice, bounded by their values on $\Z^d$), and \cref{lem:a06A:twosided,lem:a06A:domdiff} and the second bound of \cref{prop:a06A:Ba1P22} use the printed theorem only in the form just established. Hence \cref{thm:a06A:groupA} holds on $T$ with the same constants.
\end{proof}

